\section{Maximum Width of a Binary Tree}
\label{sec:trees-max-width}

\problemstatement{The width of a level is the number of positions between
its leftmost and rightmost non-null nodes, \emph{including} the null gaps
between them (as if the tree were a complete binary tree). Return the
maximum width over all levels.}

\begin{patternbox}
The phrase \emph{``including null gaps''} is the whole difficulty:
counting only present nodes is wrong. The fix is to give every node a
\textbf{positional index} as if the tree were complete -- a node at index
$i$ has children at $2i$ and $2i+1$ -- and measure each level's width as
(rightmost index $-$ leftmost index $+1$). This is BFS with indices.
\end{patternbox}

\subsection*{Understanding the problem carefully}

In a complete binary tree numbered from the root as index $0$, a node at
index $i$ places its left child at $2i+1$ and right child at $2i+2$. These
indices encode horizontal position \emph{including} the empty slots. So for
each level, the first and last real nodes carry indices whose difference,
plus one, is exactly the width that counts the gaps between them. To keep
the indices from overflowing on deep trees, subtract the level's leftmost
index from every node's index (re-basing each level to start at $0$).

\begin{ideabox}[Number Nodes as if the Tree Were Complete]
BFS level by level, pairing each node with an index (root $=0$; children
$2i+1$, $2i+2$). For each level, width $=$ (last index $-$ first index
$+1$); track the maximum. Re-base each level's indices by subtracting the
first node's index to prevent overflow.
\end{ideabox}

\subsection*{Algorithm}

\begin{enumerate}
  \item Queue holds $(node, index)$; push $(root, 0)$.
  \item For each level: let \textit{first} be the index of the level's
        first node. For every node, push $(left, 2\cdot(i-\textit{first})+1)$
        and $(right, 2\cdot(i-\textit{first})+2)$. Width $=$ last index
        $-$ first index $+1$; update the maximum.
\end{enumerate}

\subsection*{Dry run}

\dryrun{Tree: root $1$; left $2$ (left child $4$); right $3$ (right child
$5$).}

\begin{itemize}
  \item Level $0$: $(1,0)$. Width $1$.
  \item Level $1$: $(2,0),(3,1)$ (re-based). Width $1-0+1=2$.
  \item Level $2$: from $2$ push $(4,1)$; from $3$ push $(5,4)$. First
        index $1$; re-base to $(4,0),(5,3)$. Width $3-0+1=4$.
  \item Maximum width $=4$.
\end{itemize}

\subsection*{C++ implementation}

\begin{lstlisting}
#include <bits/stdc++.h>
using namespace std;

struct TreeNode {
    int val;
    TreeNode* left;
    TreeNode* right;
    TreeNode(int v) : val(v), left(nullptr), right(nullptr) {}
};

int widthOfBinaryTree(TreeNode* root) {
    if (root == nullptr) return 0;
    long long best = 0;
    queue<pair<TreeNode*, long long>> q;
    q.push({root, 0});
    while (!q.empty()) {
        int s = q.size();
        long long first = q.front().second, last = first;
        for (int i = 0; i < s; i++) {
            auto [node, idx] = q.front();
            q.pop();
            idx -= first;                 // re-base to avoid overflow
            last = idx;
            if (node->left)  q.push({node->left,  2 * idx + 1});
            if (node->right) q.push({node->right, 2 * idx + 2});
        }
        best = max(best, last + 1);
    }
    return (int)best;
}

int main() {
    TreeNode* root = new TreeNode(1);
    root->left = new TreeNode(2);
    root->right = new TreeNode(3);
    root->left->left = new TreeNode(4);
    root->right->right = new TreeNode(5);

    cout << "Max width: " << widthOfBinaryTree(root) << "\n";
    return 0;
}
\end{lstlisting}

\begin{complexitybox}
\textbf{Time Complexity.} $O(n)$ -- one BFS pass. \textbf{Space
Complexity.} $O(n)$ for the queue.
\end{complexitybox}

\begin{takeawaybox}
Positional indexing ($2i+1$, $2i+2$) is the standard way to reason about a
tree's geometry when empty slots matter. Re-basing per level is the
practical detail that keeps those indices from overflowing on a deep,
sparse tree.
\end{takeawaybox}
